\documentclass[11pt,a4paper]{article}

\usepackage[utf8]{inputenc}
\usepackage[T1]{fontenc}
\usepackage{lmodern}
\usepackage[margin=1in]{geometry}
\usepackage{amsmath,amssymb,amsthm,mathtools}
\usepackage{graphicx}
\usepackage{booktabs}
\usepackage{hyperref}
\usepackage{xcolor}
\usepackage{enumitem}

\hypersetup{
  colorlinks=true,
  linkcolor=blue!50!black,
  citecolor=blue!50!black,
  urlcolor=blue!50!black
}

\title{\bfseries Foundational and cross-disciplinary extensions \\
to the Systroph\=e framework\\
\large Five further modules (Phases 11--15)}

\author{Christian Knopp\\\texttt{cknopp@gmail.com}}
\date{2026-05-11}

\begin{document}
\maketitle

\begin{abstract}
We add five further modules to the Systroph\=e framework
\cite{Knopp2026Systrophe,Knopp2026QFTCS,Knopp2026Ext1} covering
foundational stability, observational signatures, and
cross-disciplinary applications of the Tipler / Lewis-Papapetrou
structure: (1) Cauchy horizon stability via local Lyapunov
exponents; (2) Lense-Thirring frame-dragging signatures with
comparison to Kerr; (3) optical-fiber analog (Faccio-style) of the
rotating-cylinder spacetime; (4) systematic energy-condition survey
across the $(\omega, R)$ parameter space; (5) cross-disciplinary DSI
detection in financial and biological time series, mapped to a
Systroph\=e rotation parameter $a$. Forty-six new tests across the
five modules, all passing. Code at
\texttt{github.com/Zynerji/systrophe} (v0.17.0+).
\end{abstract}

\section{Cauchy horizon stability (Phase 11)}

For each F = 0 surface in the LP supercritical exterior, we compute
the local Lyapunov exponent
\[
\lambda \sim \sqrt{|F''(r_h)|}
\]
controlling the radial geodesic divergence rate. A positive $\lambda$
indicates exponential perturbation growth -- the kinematic signature
of Poisson-Israel mass inflation \cite{PoissonIsrael1990}.

The module \texttt{cauchy\_stability.py} provides:
\begin{itemize}
\item \texttt{stability\_at\_horizon}: full descriptor at one CH;
\item \texttt{survey\_all\_cauchy\_horizons}: all CHs in range;
\item \texttt{mass\_inflation\_growth\_rate}: growth-rate estimate;
\item \texttt{chronology\_protection\_via\_instability}: full check.
\end{itemize}

For the LP supercritical exterior at $\omega R = 1$, two Cauchy
horizons appear at $r \approx 1.83$ and $r \approx 11.23$. The
Lyapunov exponents at these horizons indicate classical
\emph{instability}, supporting the Hawking chronology-protection
mechanism: small perturbations grow without bound, ultimately
preventing the closed timelike curve from forming as a stable
classical configuration.

\section{Frame-dragging signatures (Phase 12)}

The Lense-Thirring precession frequency of a stationary gyroscope at
radius $r$ on the LP background is
\[
\Omega_\text{LT}(r) = -\frac{1}{2} \frac{d}{dr}\!\left(\frac{K(r)}{r^2}\right).
\]

The module \texttt{frame\_dragging.py} computes:
\begin{itemize}
\item \texttt{lense\_thirring\_frequency}: precession rate at $r$;
\item \texttt{frame\_dragging\_pattern}: spatial profile;
\item \texttt{inertial\_drag\_zero\_locus}: where $\Omega_\text{LT} = 0$
  (frame-isolation surfaces);
\item \texttt{compare\_to\_kerr\_LT}: side-by-side comparison;
\item \texttt{total\_drag\_per\_revolution}: per-revolution precession.
\end{itemize}

Headline: for the LP supercritical exterior, $\Omega_\text{LT}(r)$
oscillates log-periodically (inherited from $K(r)$), unlike Kerr
where $\Omega_\text{LT}$ is monotonic in $r$. The cylinder thus
exhibits \emph{alternating-sign} frame-dragging zones separated by
the frame-isolation surfaces.

\section{Optical-fiber analog (Phase 13)}

Following Philbin et al.\ 2008 \cite{Philbin2008} and Faccio 2010
\cite{Faccio2010}, slow-light propagation in nonlinear optical
fibers admits an effective acoustic metric. For a pump pulse with
group velocity $v_g$ and a probe field with phase velocity $c_g$:
\[
F_\text{eff}(x) = c_g^2 - v_g(x)^2.
\]

The analog horizon is at $v_g = c_g$; the analog Hawking temperature
is $T_H = \kappa / (2\pi)$ with $\kappa = (1/2)|dF/dx|$.

The module \texttt{optical\_fiber\_analog.py} provides:
\begin{itemize}
\item \texttt{fiber\_effective\_metric}: F, K, L for the analog;
\item \texttt{fiber\_analog\_horizon}: where $v = c_g$;
\item \texttt{fiber\_analog\_hawking\_T}: predicted phonon-like T;
\item \texttt{compare\_to\_steinhauer\_2010}: comparison to fiber
  measurement of $T_H \sim 10$ $\mu$K.
\item \texttt{pulse\_collision\_radiation}: two-pulse collision case.
\end{itemize}

This provides a second experimental platform (in addition to BECs,
Phase 4) for testing the Systroph\=e Hawking predictions.

\section{Energy condition survey (Phase 14)}

The module \texttt{energy\_condition\_survey.py} runs systematic
checks of WEC, NEC, DEC, SEC across a grid of $(\omega, R)$ values.

For the classical van Stockum interior dust, all four conditions are
satisfied trivially (positive-energy dust). The interesting question
is the \emph{induced} stress-energy from QFT effects on the LP
exterior, computed via the Hadamard off-trace tensor from
\texttt{hadamard\_offtrace}.

The function \texttt{exterior\_violations\_from\_qft} reports the
WEC and NEC proxies at sample exterior points; these can be either
positive or negative depending on local curvature.

\section{DSI cross-disciplinary detection (Phase 15)}

The Tipler log-frequency $\alpha = \sqrt{4 a^2 - 1}$ is a specific
example of the broader class of discrete-scale-invariance (DSI)
signatures (Sornette 1998 \cite{Sornette1998}). Many domains exhibit
log-periodic structure:
\begin{itemize}
\item Financial crashes (Johansen-Sornette log-periodic precursors);
\item Earthquake recurrence times;
\item Biological allometry deviations;
\item Crystal-growth oscillations;
\item Critical-phenomena scaling.
\end{itemize}

The module \texttt{dsi\_crossdisciplinary.py} provides:
\begin{itemize}
\item \texttt{fit\_log\_periodic\_to\_data}: extract $\alpha$ and
  amplitude from arbitrary $(x, y)$ data;
\item \texttt{generate\_synthetic\_tipler\_signal}: known-$a$ data
  for recovery testing;
\item \texttt{financial\_crash\_log\_periodic\_signature}: apply to
  pre-crash price data;
\item \texttt{biological\_scaling\_log\_periodic}: allometric
  detrending and fit.
\end{itemize}

The fitted $\alpha$ maps to an equivalent Systroph\=e rotation
parameter via
\[
a_\text{equiv} = \sqrt{(\alpha^2 + 1) / 4}.
\]

This is a cross-disciplinary tool: any field exhibiting log-periodic
behavior can be mapped to a ``virtual cylinder'' parameter.

\section{Implications}

The five modules collectively reinforce that the Systroph\=e
framework is:
\begin{enumerate}
\item \textbf{Foundationally robust}: Cauchy horizons are
unstable in the LP supercritical regime, consistent with Hawking
chronology protection.
\item \textbf{Observationally rich}: Lense-Thirring zones plus
log-periodic optics give distinctive signatures.
\item \textbf{Experimentally accessible}: BEC vortices AND optical
fibers both serve as analog platforms.
\item \textbf{Theoretically well-grounded}: energy conditions are
satisfied classically; QFT corrections are state-dependent and
computable.
\item \textbf{Cross-disciplinarily useful}: Tipler log-frequency
maps to non-physics DSI processes via a simple parameter relation.
\end{enumerate}

\section{Reproduction}

\begin{center}
\begin{tabular}{ll}
\toprule
Phase & Source file \\
\midrule
11. Cauchy stability & \texttt{cauchy\_stability.py} \\
12. Frame-dragging & \texttt{frame\_dragging.py} \\
13. Optical fiber analog & \texttt{optical\_fiber\_analog.py} \\
14. Energy condition survey & \texttt{energy\_condition\_survey.py} \\
15. DSI cross-disciplinary & \texttt{dsi\_crossdisciplinary.py} \\
\bottomrule
\end{tabular}
\end{center}

All 46 tests pass.

\bibliographystyle{plain}
\begin{thebibliography}{9}

\bibitem{Knopp2026Systrophe}
C.~Knopp, Systrophe whitepaper I (classical GR), 2026.

\bibitem{Knopp2026QFTCS}
C.~Knopp, Systrophe whitepaper II (QFT), 2026.

\bibitem{Knopp2026Ext1}
C.~Knopp, Systrophe extensions paper (Phases 1--10), 2026.

\bibitem{PoissonIsrael1990}
E.~Poisson and W.~Israel, Phys. Rev. D 41 (1990) 1796.

\bibitem{Philbin2008}
T.~G.~Philbin et al., Science 319 (2008) 1367.

\bibitem{Faccio2010}
D.~Faccio, New J. Phys. 12 (2010) 095004.

\bibitem{Sornette1998}
D.~Sornette, Phys. Rep. 297 (1998) 239.

\end{thebibliography}

\end{document}
